\documentclass[11pt]{letter}
\usepackage[margin=1in]{geometry}
\usepackage{microtype}
\usepackage[T1]{fontenc}
\usepackage{lmodern}
\usepackage{parskip}
\signature{The Authors}
\address{}
\begin{document}
\begin{letter}{The Editors\\ \textit{[Target journal --- e.g.\ Psychological Methods / Behavior Research Methods]}}

\opening{Dear Editors,}

We submit for your consideration the manuscript \textbf{``A transdiagnostic dimensional map as
a self-calibrating measurement instrument: projection, information, and value of information
under structured missingness.''} We believe it is well suited to a quantitative-methods venue
because its contribution is methodological rather than substantive: it reframes a fitted
dimensional model of psychopathology as a \emph{measurement instrument} whose per-patient
behaviour --- how sharply each person is located, which item would sharpen them next, and how
far the reported uncertainty can be trusted --- is computable in closed form on a single fitted
object.

\textbf{The core idea.} Dimensional maps of psychopathology are typically read as fixed
point-coordinates that hide how each patient was measured. We show instead that a
missingness-native Bayesian latent-variable model, fit to each patient's observed cells only
(no imputation) under mixed likelihoods, returns for every one of $N=9{,}013$ patients a
posterior coordinate \emph{and} its uncertainty, and that four instrument properties follow
from that one object by exact identities: (i) a Woodbury-reduced expected-a-posteriori
projection that returns the prior --- never a false-precise point --- on unmeasured axes;
(ii) an information accounting showing that what localizes a patient is the Fisher information
an item carries, not its loading; (iii) a value-of-information rule, closed-form via the
matrix-determinant lemma, for designing efficient batteries; and (iv) a simulation-validated,
misspecification-stressed calibration of the reported intervals. A five-corner archetypal
simplex then summarizes the continuum as graded membership in named clinical poles.

\textbf{What is new in this version.} We have strengthened the paper specifically on the axis a
methods reviewer cares about most --- \emph{does the instrument beat the simpler tools a reader
already has?} --- by adding four head-to-head analyses. (1) Carrying the full posterior
coordinate improves out-of-sample prognosis over a naive per-axis sum-score, and the advantage
\emph{widens} as measurement grows sparse (the regime where a point estimate is least
trustworthy). (2) Ranking battery items by information rather than by loading raises mean
reliability by $0.11$ at a fixed $27$-item budget. (3) The posterior calibration degrades
gracefully --- never catastrophically --- under four controlled violations of the model's own
assumptions, converting an internal-consistency check into a robustness result. (4) We report,
with equal prominence, two honest null results (a fully patient-adaptive item order coincides
with the fixed order; and real-data confidence-triage does not improve one coarse clinical
endpoint), because each delimits precisely what the instrument does and does not buy.

We also correct and sharpen the archetype analysis: rather than defend the simplex on
reconstruction fidelity --- which we now show is the \emph{wrong} yardstick, since the map's
near-isotropic variance lets even a random five-dimensional subspace retain more variance than
the archetypes --- we make the case on what the simplex uniquely provides over principal
components: convex, non-negative, per-patient membership in five nameable clinical poles, over a
population we confirm has no natural clusters.

\textbf{Relationship to a companion paper.} This manuscript is the methodological half of a
two-paper pair. A companion clinical paper (in preparation / under separate submission) reports
the substantive findings from the same cohort --- the biology-general-factor separation, the
durable axis, and the prognostic gradient across archetypes. The two papers share the fitted
model but are disjoint in claims: the present paper makes no clinical discovery claim, and the
companion makes no methodological one. We are happy to provide the companion manuscript to the
editors on request to confirm the absence of overlap.

\textbf{Declarations.} The work uses no external data beyond the cohort described; all analysis
code and the fitted-model scoring scripts are released with the paper. The manuscript is not
under consideration elsewhere, and all authors approve this submission. We have no competing
interests to declare.

We would be glad to suggest reviewers with expertise in item-response and computerized adaptive
testing, latent-variable models under missing data, and Bayesian measurement calibration, and we
thank you for considering the manuscript.

\closing{With best regards,}

\end{letter}
\end{document}